\section{Related work}
\label{sec:related}

\subsection{Coercions and subtyping}

Implicit coercions have a long theory. \citet{breazu1991} interpret inheritance as
implicit coercion and prove \emph{coherence}: the meaning of a term does not depend on
which typing derivation, and so which insertion of coercions, is chosen.
\citet{luo1999} develops coercive subtyping for dependent type theories, in which
coercions are an abbreviation mechanism and coherence, the uniqueness of the coercion
between two types, is the central requirement. Coherence answers the question
\emph{which} map is inserted. The question raised in Section~\ref{sec:casts} is the next
one: given the map, \emph{what} it preserves. The two are independent, since the unique
cast from $\N$ to $\Z$ still fails to commute with truncated subtraction.

Within Lean, the simplification of casts is described by \citet{lewis2020}, whose
classification of lemmas into move, elim and squash is used in
Section~\ref{sec:casts}, and the design of bundled morphisms and their classes by
\citet{baanen2022}. This paper adds nothing to that machinery. It reads it as an
instance of a more general practice.

The requirement that truth be invariant under a change of notation is the satisfaction
condition of the theory of institutions \citep{goguen1992}, where it is an axiom on
translations between signatures. Predicate locality asks for the same invariance, but as
an obligation to be discharged for each predicate that is reused across a boundary, and
with the expectation that it will sometimes fail.

\subsection{Refinement and the preservation of properties}

That refinement preserves properties of single traces and not, in general, properties of
sets of traces is the subject of \citet{clarkson2010}; the distinction between safety and
liveness on which Section~\ref{sec:deontic} relies is due to \citet{alpern1985}. The
theory of secure compilation makes the preservation list explicit. \citet{abate2019}
define a family of criteria for compilers according to the class of properties that is
preserved against adversarial contexts (trace properties, safety properties,
hyperproperties and relational hyperproperties) and establish how the criteria relate;
\citet{patrignani2019} survey the earlier criterion of full abstraction. These criteria
are the most developed existing counterpart of the second component of a typed crossing:
they type a compiler by what it preserves.

The present paper is less precise than that work and differs in scope. It adds the first
component, the kind of the crossing, and it includes among the things a crossing may fail
to carry some that are not properties of traces at all: duration
(Section~\ref{sec:duration}), provenance, and the standing of a step under a norm.

\subsection{Faithfulness, causal abstraction and monitoring}

\citet{jacovi2020} separate the faithfulness of an explanation from its plausibility and
argue for a graded notion of faithfulness in place of a binary one. The model in
Section~\ref{sec:kernel} uses a binary relation, which is the simplest setting in which
the point about checks on text can be made; a graded notion would need a different
formal treatment. Empirical work on written reasoning includes \citet{turpin2023},
\citet{lanham2023} and \citet{chen2025}. \citet{barez2025} argue from this evidence that
a chain of thought should not be treated as an explanation without causal validation,
which is the conclusion of Section~\ref{sec:reasoning} reached from the empirical side.

The formal theory of the crossing between an explanation and a mechanism is causal
abstraction \citep{geiger2025}: a high-level causal model is an abstraction of a
low-level one when interventions on the two correspond, and this is tested by
intervening on the network. Proposition~\ref{prop:char} is far weaker than that theory
and compatible with it. Its one point is conditional: where faithfulness differs
between mechanisms that produce the same description, an exact judgement requires
information beyond the description. Interventions on the network, as used in causal
abstraction, are one way to obtain such information.

\citet{korbak2025} argue that written reasoning is nevertheless a valuable signal for
monitoring models, and a fragile one, since training methods or architectures may make
it less legible. That position is consistent with this paper. A monitor that reads a
trace obtains evidence, by an evidential edge, about what the model may do. The results
here concern a different question, whether a check on the trace can decide that the
trace describes the computation, and Remark~\ref{rem:scope}(b) keeps the two apart.

\subsection{Levels and edge types}

The distinction between constitutive and causal relations between levels is taken from
the mechanistic account of \citet{craver2007}, and the intervention test for causal
edges from \citet{woodward2003}. The edge taxonomy, the five structures and the axioms
are those of LEM \citep{kreinzlem}, which discusses its own relation to the literature
on levels, ontology and risk modelling.
